\documentclass[12pt]{article}
\usepackage{amsmath}
\usepackage{amssymb}
\usepackage{amsfonts}
\usepackage{amsthm}
\usepackage{bm}
\usepackage{array}
\usepackage[margin=1in]{geometry}
\usepackage[colorlinks=true,linkcolor=blue,citecolor=blue,urlcolor=blue]{hyperref}

\numberwithin{equation}{section}

\newtheorem{theorem}{Theorem}
\newtheorem{proposition}{Proposition}
\newtheorem{lemma}{Lemma}
\newtheorem{corollary}{Corollary}
\newtheorem{assumption}{Assumption}
\theoremstyle{definition}
\newtheorem{definition}{Definition}
\newtheorem{remark}{Remark}

\newcommand{\ind}{\mathbf 1}
\newcommand{\R}{\mathbb R}
\newcommand{\E}{\mathbb E}
\newcommand{\Pp}{\mathbb P}
\newcommand{\F}{\mathcal F}
\newcommand{\D}{\mathcal D}
\newcommand{\op}{o_p}
\newcommand{\Op}{O_p}

\begin{document}

\begin{center}
{\Large\bf Bai--Perron Theory Workbook: Conditional-Design Rebuild}\\[6pt]
{\large Profile Sieve Empirical Likelihood with Nuisance Breaks}\\[6pt]
{\normalsize One-break estimated-partition Wilks, unknown-order selection, and proof map}
\end{center}

\tableofcontents
\newpage

\section{Primitive Setup and Notation}

For a scalar array \(\bm a=(a_1,\ldots,a_T)'\), write
\[
    \|\bm a\|_T^2=T^{-1}\sum_{t=1}^T a_t^2,
    \qquad
    \|\bm a\|_\infty=\max_{t\leq T}|a_t|.
\]
All \(O_p(\cdot)\) and \(o_p(\cdot)\) statements for vectors and matrices use Euclidean and spectral norm unless another norm is displayed.

The maintained null model is
\begin{equation}
    Y_t
    =
    \beta_0X_{t-1}
    +
    \sum_{j=0}^{q_0}
    m_j(W_{t-1})\ind\{k_{j,0}<t\leq k_{j+1,0}\}
    +
    U_t,
    \label{eq:primitive-model}
\end{equation}
where
\[
    0=k_{0,0}<k_{1,0}<\cdots<k_{q_0,0}<k_{q_0+1,0}=T.
\]
Throughout the known-partition theory, \(q_0\) is fixed.  The main
estimated-partition theorem later specializes to one unknown nuisance break,
\(q_0=1\).  Candidate searches use a trimming constant \(\epsilon\), while
the true partition has slack \(\epsilon_0>\epsilon\):
\[
    \min_{0\leq j\leq q_0}(k_{j+1,0}-k_{j,0})\geq \epsilon_0 T,
    \qquad
    (\bar q+1)\epsilon<1.
\]
The coefficient \(\beta_0\) is stable. All structural breaks are in the nuisance component.

For a candidate partition \(\Lambda=(k_1,\ldots,k_q)\), let \(k_0=0\), \(k_{q+1}=T\), and
\[
    I_j(\Lambda)=\{t:k_j<t\leq k_{j+1}\},
    \qquad
    j_\Lambda(t)=j\ \text{if }t\in I_j(\Lambda).
\]
For the search trimming constant, define the admissible candidate set
\[
    \mathcal L_q(\epsilon)
    =
    \left\{\Lambda=(k_1,\ldots,k_q):
    0=k_0<k_1<\cdots<k_q<k_{q+1}=T,\ 
    \min_{0\leq j\leq q}(k_{j+1}-k_j)\geq \epsilon T
    \right\}.
\]
Let
\[
    \mu_t=m_{j_{\Lambda_0}(t)}(W_{t-1}),
    \qquad
    \bm\mu=(\mu_1,\ldots,\mu_T)'.
\]

Let \(\bm P\) be the \(T\times K\) sieve matrix with \(t\)-th row
\[
    \bm P_K(W_{t-1})'=(p_1(W_{t-1}),\ldots,p_K(W_{t-1})).
\]
For a partition \(\Lambda\), let \(\bm D_j(\Lambda)\) be the \(T\times T\) diagonal selector matrix with \(t\)-th diagonal entry \(\ind\{t\in I_j(\Lambda)\}\), \(j=0,\ldots,q\). Define
\begin{equation}
    \bm Q_{K,\Lambda}
    =
    [\bm D_0(\Lambda)\bm P,\ldots,\bm D_q(\Lambda)\bm P],
    \qquad
    \bm M_{K,\Lambda}
    =
    \bm I_T-\bm Q_{K,\Lambda}
    (\bm Q_{K,\Lambda}'\bm Q_{K,\Lambda})^\dagger
    \bm Q_{K,\Lambda}'.
    \label{eq:primitive-qm}
\end{equation}
The Moore--Penrose inverse makes \(\bm M_{K,\Lambda}\) algebraically defined even when \(\bm Q_{K,\Lambda}'\bm Q_{K,\Lambda}\) is singular. Such singularity can arise if a candidate segment is too short relative to \(K\) or if the selected basis columns are collinear on that segment; the projection remains defined, but the asymptotic regularity used below may fail. The theory therefore relies on minimum-spaced admissible partitions and Gram-stability/full-rank conditions on the relevant candidate sets.

For each \((\beta,\Lambda)\), define
\begin{align}
    \widehat{\bm u}(\beta,\Lambda)
    &=
    \bm M_{K,\Lambda}(\bm Y-\beta\bm X_{lag}),
    \\
    \bm w^c(\Lambda)
    &=
    \bm M_{K,\Lambda}\bm w,
    \\
    \widehat Z_t(\beta,\Lambda)
    &=
    \widehat u_t(\beta,\Lambda)w_t^c(\Lambda).
\end{align}
Let
\[
    S_T(\Lambda)
    =
    T^{-1/2}\sum_{t=1}^T\widehat Z_t(\beta_0,\Lambda),
    \qquad
    V_T(\Lambda)
    =
    T^{-1}\sum_{t=1}^T\widehat Z_t(\beta_0,\Lambda)^2.
\]

For a scalar score array \(Z_1,\ldots,Z_T\) satisfying the convex-hull
condition, let \(\lambda_T\) be the unique solution of
\[
    \sum_{t=1}^T\frac{Z_t}{1+\lambda_T Z_t}=0,
    \qquad
    1+\lambda_T Z_t>0\quad(t\leq T),
\]
and define the twice negative empirical-likelihood log ratio by
\begin{equation}
    \ell_T(Z_1,\ldots,Z_T)
    =2\sum_{t=1}^T\log(1+\lambda_TZ_t).
    \label{eq:sol-el-definition}
\end{equation}
In particular,
\[
    \ell_T(\beta,\Lambda)
    =\ell_T\{\widehat Z_1(\beta,\Lambda),\ldots,
                    \widehat Z_T(\beta,\Lambda)\}.
\]

For an interval \([s,e]\), define the null-imposed segment cost
\[
    \mathcal C_K(s,e;\beta)
    =
    \min_{\bm c\in\R^K}
    \sum_{t=s}^e
    \{Y_t-\beta X_{t-1}-\bm P_K(W_{t-1})'\bm c\}^2.
\]
For \(\Lambda=(k_1,\ldots,k_q)\), define
\[
    RSS_K(\beta,\Lambda)
    =
    \sum_{j=0}^q
    \mathcal C_K(k_j+1,k_{j+1};\beta).
\]
All minimizers are measurable selections with a fixed deterministic tie-breaking rule. For known order,
\[
    \widehat\Lambda_q(\beta)
    =
    \arg\min_{\Lambda\in\mathcal L_q(\epsilon)}
    RSS_K(\beta,\Lambda),
    \qquad
    \widetilde\Lambda=\widehat\Lambda_{q_0}(\beta_0).
\]
For unknown order,
\[
    \widehat q(\beta)
    =
    \arg\min_{0\leq q\leq \bar q}
    \{RSS_K(\beta,\widehat\Lambda_q(\beta))+\operatorname{pen}_T(q,K)\},
    \qquad
    \widehat\Lambda(\beta)=\widehat\Lambda_{\widehat q(\beta)}(\beta).
\]
Although the Moore--Penrose inverse makes \(\bm M_{K,\Lambda}\) algebraically defined for arbitrary \(\Lambda\), all theoretical searches are restricted to minimum-spaced admissible partitions \(\mathcal L_q(\epsilon)\); thus no candidate segment is allowed to be asymptotically shorter than the sieve dimension \(K\), and the Gram matrices used in the proofs are nonsingular on the relevant candidate set.

The sieve dimension is allowed to depend on the sample size and is written \(K=K_T\). Throughout, \(K_T\to\infty\), subject to the rate restrictions in Assumptions \ref{ass:prim-bounded}--\ref{ass:prim-sieve} and the local Gram regularity conditions below. The basis is normalized so that its first component is constant, say \(p_1(w)\equiv1\), allowing regime-specific intercept and level shifts to be represented by the block sieve.

For local-power calculations, distinguish the null value used in the statistic from the data-generating slope. The local alternative used in the rebuilt theorem is
\[
    \beta_T=\beta_0+\frac{h}{\sqrt T},
    \qquad h\in\R,
\]
with data generated under \(P_{\beta_T}\) while the profile empirical likelihood is evaluated at \(\beta_0\). Equivalently,
\[
    Y_t-\beta_0X_{t-1}=\mu_t+U_t+\frac{h}{\sqrt T}X_{t-1},
\]
with the same nuisance break partition \(\Lambda_0\). Section \ref{sec:local-alternatives} records the corresponding drift conditions.

\section{Primitive Assumptions}

\begin{assumption}[Conditional-on-design stochastic environment]
\label{ass:prim-stoch}
Let
\[
    \D_T=\sigma\{X_{t-1},W_{t-1}:1\leq t\leq T\}.
\]
Conditional on \(\D_T\), the variables \(U_{1,T},\ldots,U_{T,T}\) are
independent and satisfy
\[
    \E(U_{t,T}\mid\D_T)=0,
    \qquad
    0<\underline\sigma^2\leq
    \sigma_{t,T}^2:=\E(U_{t,T}^2\mid\D_T)
    \leq\overline\sigma^2<\infty.
\]
Moreover, for every \(s\in\R\),
\[
    \E\{\exp(sU_{t,T})\mid\D_T\}
    \leq \exp(Cs^2/2)
\]
almost surely, uniformly in \(t\) and \(T\).  The subscript \(T\) is suppressed
below when no ambiguity can arise.
\end{assumption}

\begin{assumption}[Design-measurable score regularity]
\label{ass:prim-score-admissible}
For every admissible partition \(\Lambda\), the matrices \(\bm Q_{K,\Lambda}\),
\(\bm M_{K,\Lambda}\), \(\bm\Pi_{K,\Lambda}=\bm I_T-\bm M_{K,\Lambda}\), and
\(\bm v_\Lambda=\bm M_{K,\Lambda}\bm w\) are \(\D_T\)-measurable.  Hence the
full-sample residualized instrument creates no predictability problem: all
oracle score calculations are conditional-on-design calculations.  For the true
partition, write \(\bm v_0=\bm v_{\Lambda_0}\).
\end{assumption}

\begin{assumption}[Bounded weights, nuisance functions, and basis envelope]
\label{ass:prim-bounded}
The sieve dimension satisfies \(K=K_T\to\infty\). The first basis element is constant, \(p_1(w)\equiv1\). There are constants \(C_w,C_m<\infty\) and an envelope \(\zeta_K\) such that
\[
    \max_{t\leq T}|w_t|\leq C_w,
    \qquad
    \max_{0\leq j\leq q_0}\sup_x|m_j(x)|\leq C_m,
    \qquad
    \max_{t\leq T}\|\bm P_K(W_{t-1})\|\leq \zeta_K
\]
with probability tending to one.
\end{assumption}

\begin{assumption}[Spacing, smoothness, and break detectability]
\label{ass:prim-breaks}
The known-partition theory allows fixed \(q_0\).  The estimated-partition and
unknown-order Wilks theorems in the main text impose \(q_0=1\).  The true
segments obey the strict spacing condition in the setup with
\(\epsilon_0>\epsilon\), and searches are over \(\mathcal L_q(\epsilon)\) for
\(0\leq q\leq\bar q\).  Each \(m_j\) belongs to a common smoothness class
\(\mathcal H^s\).  If \(q_0\geq1\), define the design jump size by
\[
    \Delta_T=\min_{0\leq j<q_0}
    \left\{T^{-1}\sum_{t=1}^T
    [m_{j+1}(W_{t-1})-m_j(W_{t-1})]^2\right\}^{1/2}.
\]
For the one-break main theorem, this reduces to the empirical jump norm between
\(m_0\) and \(m_1\), and \(\Delta_T>0\).
\end{assumption}

\begin{assumption}[Global and local Gram stability]
\label{ass:prim-gram}
Uniformly over every interval \(J\) that can appear as a segment in an
admissible search with \(|J|\geq\epsilon T\),
\[
    c_G I_K\preceq |J|^{-1}\bm P_J'\bm P_J\preceq C_G I_K
\]
with probability tending to one.  The same bounds hold on all one-break local
segments satisfying \(|k-k_0|\leq Cr_T\).  This condition is stronger than
algebraic Moore--Penrose existence: it supplies the inverse and leverage bounds
used by the break bridge and by the profile-RSS search.
\end{assumption}

\begin{assumption}[Sieve approximation and projection bias]
\label{ass:prim-sieve}
For each regime \(j\), there is a coefficient vector \(\bm c_{j,K}\) such that
\[
    m_j(W_{t-1})
    =
    \bm P_K(W_{t-1})'\bm c_{j,K}+r_{j,K,t},
    \qquad
    \max_j\|\bm r_{j,K}\|_T=O_p(a_K),
    \qquad
    \max_j\|\bm r_{j,K}\|_\infty=O_p(a_K).
\]
Let \(\bm r_0\) be the stacked true-partition approximation error. The rates satisfy
\[
    T^{-1/2}\bm r_0'\bm M_{K,\Lambda_0}\bm w=o_p(1),
    \qquad
    a_K(1+\zeta_K)\to0,
    \qquad
    a_K(1+\zeta_K)^2=o(\sqrt T).
\]
Also,
\[
    \frac{(q_0+1)K}{\sqrt T}\to0,
    \qquad
    \zeta_K\sqrt{\frac{(q_0+1)K}{T}}\to0.
\]
\end{assumption}

\begin{assumption}[One-break localization input for the bridge]
\label{ass:prim-localization-target}
In the one-break bridge, \(q_0=1\), \(\Lambda_0=(k_0)\), and
\(\widetilde\Lambda=(\widetilde k)\).  The bridge uses
\[
    |\widetilde k-k_0|=O_p(r_T),
    \qquad
    r_T=o(\sqrt T).
\]
Section \ref{sec:break-localization} derives this rate from the
conditional-design profile-sieve segmentation conditions with
\(r_T=R_T/\Delta_T^2\).
\end{assumption}

\begin{assumption}[Empirical-likelihood convex hull]
\label{ass:prim-convex}
For the score array under consideration,
\[
    \Pp\{\min_{t\leq T}Z_{t,T}<0<\max_{t\leq T}Z_{t,T}\}\to1.
\]
When used for the true partition, \(Z_{t,T}=\widehat Z_t(\beta_0,\Lambda_0)\). When used for the estimated partition, \(Z_{t,T}=\widehat Z_t(\beta_0,\widetilde\Lambda)\).
\end{assumption}


\section{Conditional-Design Completing Conditions}
\label{sec:rigour-audit}

The rebuilt proof uses one stochastic environment throughout.  Conditional on
\(\D_T\), the sieve design, all projection matrices, residualized weights, and
break estimators are fixed.  The remaining conditions below are rate and
identification conditions; they do not reintroduce a predictable-score route.

Write
\[
    d_K=(q_0+1)K,
    \qquad b_K=1+\zeta_K.
\]

\begin{assumption}[Oracle envelope completion]
\label{ass:sol-oracle-envelope}
For the true partition,
\[
    \max_{t\leq T}|v_{t,0}|=O_p(b_K),
    \qquad
    \frac{b_K^2\log T}{T}\longrightarrow0.
\]
The first display follows from the leverage bound under Assumptions
\ref{ass:prim-bounded} and \ref{ass:prim-gram}; it is stated here to make the
known-partition Wilks theorem self-contained.  The logarithmic rate is the
conditional sub-Gaussian maximum-score rate used below.
\end{assumption}

\begin{assumption}[Local sieve-bias stability]
\label{ass:sol-local-bias}
In the one-break section, for every fixed \(C<\infty\),
\begin{equation}
    \sup_{|k-k_0|\leq Cr_T}
    T^{-1/2}\left|
      \bm r_0'\{\bm M_{K,k}-\bm M_{K,k_0}\}\bm w
    \right|=o_p(1).
    \label{eq:sol-local-bias}
\end{equation}
A sufficient primitive rate, proved by the coefficient-update argument in
Lemma \ref{lem:wb-bridge-mu}, is
\begin{equation}
    \frac{a_K b_K(1+\zeta_K^2)r_T}{\sqrt T}\longrightarrow0.
    \label{eq:sol-local-bias-sufficient}
\end{equation}
\end{assumption}

\begin{assumption}[Local full-rank series regularity]
\label{ass:sol-local-series}
In the one-break section, every segment Gram matrix associated with a
partition satisfying \(|k-k_0|\leq Cr_T\) is nonsingular, and its eigenvalues
are between \(cT\) and \(CT\), uniformly with probability tending to one.
Also
\begin{equation}
    \zeta_K^2=O(K).
    \label{eq:sol-regular-envelope}
\end{equation}
The last display is the usual normalized-series envelope condition.  It is
used to simplify the denominator bridge rate; without it the larger order
\(b_K^2\zeta_K^2r_T/T\) should be retained.
\end{assumption}

\begin{assumption}[Primitive conditional-design segmentation route]
\label{ass:sol-localization}
This is the primitive route used to discharge the profile-sieve break search in
the one-break main theorem.  Let \(\bm\mu_K=\bm Q_{K,\Lambda_0}\bm c_0\) be the
sieve part of the true nuisance mean and define the deterministic conditional
profile loss
\[
    D_T(\Lambda)=\|\bm M_{K,\Lambda}\bm\mu_K\|^2.
\]
At the true partition, \(D_T(\Lambda_0)=0\).  For one-break candidates,
\[
    D_T(k)\geq c_\Delta |k-k_0|\Delta_T^2.
\]
For underfitted candidates \(q<q_0\),
\[
    \inf_{\Lambda\in\mathcal L_q(\epsilon)}D_T(\Lambda)
    \geq c_\Delta T\Delta_T^2.
\]
The approximation error obeys the sup-norm rate in Assumption
\ref{ass:prim-sieve}; write
\[
    R_T=K+\log T+Ta_K^2,
    \qquad
    r_T=R_T/\Delta_T^2.
\]
The rates satisfy
\[
    R_T=o(T\Delta_T^2),
    \qquad
    r_T=o(\sqrt T).
\]
Under Assumptions \ref{ass:prim-stoch} and \ref{ass:prim-gram}, conditional
sub-Gaussian concentration for the finite admissible search then gives the
profile-RSS localization and the overfit/underfit RSS gaps recorded in
Assumption \ref{ass:bp-profile-sieve}.
\end{assumption}

\begin{remark}[Why the rebuild uses conditional design]
The former martingale version needed a predictable approximation to the
full-sample residualized instrument.  The conditional-design route avoids that
problem: \(\bm M_{K,\Lambda}\), \(\bm v_\Lambda\), and the RSS minimizer are
all \(\D_T\)-measurable, so stochastic arguments are conditional concentration
and conditional triangular-array CLT arguments.  The cost is a stronger
exogeneity assumption and, in the main estimated-break Wilks theorem, the
explicit one-break scope \(q_0=1\).
\end{remark}
\begin{lemma}[Scalar empirical-likelihood expansion]
\label{lem:sol-scalar-el}
Let \(Z_{1,T},\ldots,Z_{T,T}\) satisfy the convex-hull condition and suppose
\[
    A_T=T^{-1/2}\sum_{t=1}^TZ_{t,T}=O_p(1),
    \qquad
    B_T=T^{-1}\sum_{t=1}^TZ_{t,T}^2\geq c+o_p(1)
\]
for some \(c>0\), and
\[
    \max_{t\leq T}|Z_{t,T}|=o_p(\sqrt T).
\]
Then
\begin{equation}
    \ell_T(Z_{1,T},\ldots,Z_{T,T})
    =\frac{A_T^2}{B_T}+o_p(1).
    \label{eq:sol-scalar-el-expansion}
\end{equation}
\end{lemma}

\begin{proof}
Put
\(S_1=\sum_tZ_{t,T}\), \(S_2=\sum_tZ_{t,T}^2\), and
\(M_T=\max_t|Z_{t,T}|\).  The convex-hull condition gives a unique multiplier
\(\lambda_T\) in the interval on which every \(1+\lambda_TZ_{t,T}\) is
positive.  Its score equation implies the exact identity
\[
    S_1
    =\lambda_T\sum_{t=1}^T
      \frac{Z_{t,T}^2}{1+\lambda_TZ_{t,T}}.
\]
Consequently,
\[
    |S_1|
    \geq
    \frac{|\lambda_T|S_2}{1+|\lambda_T|M_T}.
\]
After rearrangement,
\[
    |\lambda_T|\{S_2-|S_1|M_T\}\leq |S_1|.
\]
Now \(|S_1|=O_p(\sqrt T)\), \(S_2\geq(c+o_p(1))T\), and
\(|S_1|M_T=o_p(T)\).  Hence
\[
    \lambda_T=O_p(T^{-1/2}),
    \qquad
    |\lambda_T|M_T=o_p(1).
\]
Using \((1+x)^{-1}=1-x+x^2/(1+x)\) in the score equation gives
\[
    0=S_1-\lambda_TS_2+R_{1T},
    \qquad
    |R_{1T}|
    \leq
    \frac{\lambda_T^2M_TS_2}{1-|\lambda_T|M_T}.
\]
Thus
\[
    \lambda_T=\frac{S_1}{S_2}+o_p(T^{-1/2}).
\]
Finally, Taylor expansion of the logarithm, uniformly because
\(|\lambda_T|M_T=o_p(1)\), yields
\[
    2\sum_t\log(1+\lambda_TZ_{t,T})
    =2\lambda_TS_1-\lambda_T^2S_2+R_{2T},
\]
where
\[
    |R_{2T}|
    \leq C|\lambda_T|^3M_TS_2=o_p(1).
\]
Substituting the expansion for \(\lambda_T\) gives
\[
    2\sum_t\log(1+\lambda_TZ_{t,T})
    =\frac{S_1^2}{S_2}+o_p(1)
    =\frac{A_T^2}{B_T}+o_p(1),
\]
which proves \eqref{eq:sol-scalar-el-expansion}.
\end{proof}

\section{Deterministic Projection Algebra}

\begin{lemma}[Block projection identities]
\label{lem:wb-projection}
For every partition \(\Lambda\),
\[
    \bm Q_{K,\Lambda}'\bm M_{K,\Lambda}=\bm0,
    \qquad
    \bm M_{K,\Lambda}'=\bm M_{K,\Lambda},
    \qquad
    \bm M_{K,\Lambda}^2=\bm M_{K,\Lambda}.
\]
Consequently, for every conformable \(\bm c\),
\[
    (\bm Q_{K,\Lambda}\bm c)'\bm M_{K,\Lambda}\bm w=0.
\]
\end{lemma}


\begin{proof}[Detailed solution]
Abbreviate \(\bm Q=\bm Q_{K,\Lambda}\).  Take a thin singular-value
decomposition
\[
    \bm Q=\bm U_r\bm\Sigma_r\bm V_r',
\]
where \(r=\operatorname{rank}(\bm Q)\), the columns of \(\bm U_r\) and
\(\bm V_r\) are orthonormal, and \(\bm\Sigma_r\) is diagonal with strictly
positive entries.  Then
\[
    \bm Q'\bm Q
    =\bm V_r\bm\Sigma_r^2\bm V_r',
    \qquad
    (\bm Q'\bm Q)^\dagger
    =\bm V_r\bm\Sigma_r^{-2}\bm V_r'.
\]
Therefore
\[
    \bm\Pi_{K,\Lambda}
    =\bm Q(\bm Q'\bm Q)^\dagger\bm Q'
    =\bm U_r\bm U_r'.
\]
This is the orthogonal projector onto \(\operatorname{col}(\bm Q)\).  In
particular,
\[
    \bm\Pi_{K,\Lambda}'=\bm\Pi_{K,\Lambda},
    \qquad
    \bm\Pi_{K,\Lambda}^2=\bm\Pi_{K,\Lambda},
    \qquad
    \bm\Pi_{K,\Lambda}\bm Q=\bm Q.
\]
Since \(\bm M_{K,\Lambda}=\bm I_T-\bm\Pi_{K,\Lambda}\),
\[
    \bm M_{K,\Lambda}'=\bm M_{K,\Lambda}
\]
and
\[
    \bm M_{K,\Lambda}^2
    =(\bm I_T-\bm\Pi_{K,\Lambda})^2
    =\bm I_T-2\bm\Pi_{K,\Lambda}+\bm\Pi_{K,\Lambda}^2
    =\bm M_{K,\Lambda}.
\]
Also
\[
    \bm Q'\bm M_{K,\Lambda}
    =\{\bm M_{K,\Lambda}\bm Q\}'
    =\bm0'.
\]
For conformable \(\bm c\) and \(\bm w\), it follows immediately that
\[
    (\bm Q\bm c)'\bm M_{K,\Lambda}\bm w
    =\bm c'\bm Q'\bm M_{K,\Lambda}\bm w=0.
\]
\end{proof}


\begin{lemma}[Uniform fitted-value and leverage bounds]
\label{lem:wb-leverage}
Under Assumptions \ref{ass:prim-bounded} and \ref{ass:prim-gram}, uniformly over
\[
    \mathcal N_T(r_T)=
    \{\Lambda: q=q_0,
      \max_{1\leq j\leq q_0}|k_j-k_{j,0}|\leq Cr_T\},
\]
\[
    \|\bm M_{K,\Lambda}\bm w\|_\infty=O_p(1+\zeta_K),
    \qquad
    \max_{t\leq T}\{\bm Q_{K,\Lambda}(\bm Q_{K,\Lambda}'\bm Q_{K,\Lambda})^\dagger\bm Q_{K,\Lambda}'\}_{tt}
    =
    O_p(\zeta_K^2K/T).
\]
\end{lemma}


\begin{proof}[Detailed solution]
Write \(\bm Q=\bm Q_{K,\Lambda}\), \(\bm G=\bm Q'\bm Q\), and let
\(\bm q_t'\) denote row \(t\) of \(\bm Q\).  Because exactly one regime block
is active in each row,
\[
    \|\bm q_t\|=\|\bm P_K(W_{t-1})\|\leq\zeta_K.
\]
On the uniform Gram event in Assumption \ref{ass:prim-gram},
\[
    \|\bm G^\dagger\|\leq \frac{C}{T},
    \qquad
    \|\bm Q\|^2=\lambda_{\max}(\bm Q'\bm Q)\leq CT.
\]
Since \(\|\bm w\|\leq C_w\sqrt T\), the least-squares coefficient vector
\[
    \widehat{\bm c}_\Lambda=\bm G^\dagger\bm Q'\bm w
\]
satisfies
\[
    \|\widehat{\bm c}_\Lambda\|
    \leq \frac{C}{T}\,\|\bm Q\|\,\|\bm w\|
    \leq C.
\]
Thus, uniformly in \(t\) and \(\Lambda\),
\[
    |(\bm\Pi_{K,\Lambda}\bm w)_t|
    =|\bm q_t'\widehat{\bm c}_\Lambda|
    \leq C\zeta_K,
\]
and hence
\[
    \|\bm M_{K,\Lambda}\bm w\|_\infty
    \leq \|\bm w\|_\infty+
          \|\bm\Pi_{K,\Lambda}\bm w\|_\infty
    =O_p(1+\zeta_K).
\]
For the leverage,
\[
    (\bm\Pi_{K,\Lambda})_{tt}
    =\bm q_t'\bm G^\dagger\bm q_t
    \leq \|\bm G^\dagger\|\,\|\bm q_t\|^2
    \leq C\frac{\zeta_K^2}{T}.
\]
This is sharper than the displayed target and therefore implies
\(O_p(\zeta_K^2K/T)\), because \(K\geq1\).  All inequalities hold on a single
uniform Gram event whose probability tends to one.
\end{proof}


\section{Known-Partition Conditional-Design Wilks}

Set
\[
    \bm v_0=\bm M_{K,\Lambda_0}\bm w,
    \qquad
    A_T^{(q)}=T^{-1/2}\sum_{t=1}^TU_tv_{t,0},
    \qquad
    V_T^{(q)}=T^{-1}\sum_{t=1}^TU_t^2v_{t,0}^2.
\]

\begin{lemma}[Oracle conditional-design CLT]
\label{lem:wb-oracle-clt}
Under Assumptions \ref{ass:prim-stoch}--\ref{ass:prim-gram} and
\ref{ass:sol-oracle-envelope}, suppose
\[
    \eta_{T,q}^2:=T^{-1}\sum_{t=1}^T \sigma_t^2v_{t,0}^2
    \xrightarrow{p}\eta_q^2,
    \qquad
    \Pp(\eta_q\geq \sqrt{c_\eta})=1
\]
for some \(c_\eta>0\). Then
\[
    A_T^{(q)}\xrightarrow{\mathrm{st}}\eta_qZ,
\]
where \(Z\sim N(0,1)\) is independent of the limiting sigma-field generated by
\(\eta_q\).
\end{lemma}

\begin{proof}[Detailed solution]
Condition on \(\D_T\).  The summands
\(T^{-1/2}U_tv_{t,0}\) are independent, centered, and have conditional
variance \(T^{-1}\sigma_t^2v_{t,0}^2\).  The conditional variance sum is
\(\eta_{T,q}^2\).  The conditional sub-Gaussian envelope and
\(\max_t|v_{t,0}|=O_p(b_K)\), together with
\(b_K/\sqrt T\to0\), imply the Lindeberg condition conditionally on \(\D_T\):
for every \(\varepsilon>0\),
\[
    \sum_{t=1}^T
    \E\left[ T^{-1}U_t^2v_{t,0}^2
    \ind\{|U_tv_{t,0}|>\varepsilon\sqrt T\}\mid\D_T\right]
    =o_p(1).
\]
Therefore the conditional triangular-array CLT gives
\[
    \E\{\exp(iuA_T^{(q)})\mid\D_T\}
    -\exp(-u^2\eta_{T,q}^2/2)\xrightarrow{p}0
\]
for every fixed \(u\).  Since \(\eta_{T,q}^2\to_p\eta_q^2\), the usual
conditional-characteristic-function argument yields stable convergence to
\(\eta_qZ\).  The lower bound on \(\eta_q\) is recorded for the later
studentization step.
\end{proof}

\begin{lemma}[Oracle denominator LLN]
\label{lem:wb-oracle-lln}
Under Assumptions \ref{ass:prim-stoch}--\ref{ass:prim-gram} and
\ref{ass:sol-oracle-envelope}, suppose
\[
    T^{-1}\sum_{t=1}^T \sigma_t^2v_{t,0}^2\xrightarrow{p}\eta_q^2.
\]
Then
\[
    V_T^{(q)}-\eta_q^2=o_p(1).
\]
\end{lemma}

\begin{proof}[Detailed solution]
Conditional on \(\D_T\), the variables
\((U_t^2-\sigma_t^2)v_{t,0}^2\) are independent and centered.  Conditional
sub-Gaussianity gives a uniform fourth-moment bound, hence
\[
    \operatorname{Var}\left(
    T^{-1}\sum_{t=1}^T (U_t^2-\sigma_t^2)v_{t,0}^2
    \mid\D_T\right)
    \leq CT^{-2}\sum_{t=1}^Tv_{t,0}^4.
\]
Using \(\max_t|v_{t,0}|=O_p(b_K)\) and
\(T^{-1}\sum_t v_{t,0}^2=O_p(1)\), the right side is
\(O_p(b_K^2/T)=o_p(1)\).  Thus
\[
    T^{-1}\sum_{t=1}^T U_t^2v_{t,0}^2
    -T^{-1}\sum_{t=1}^T \sigma_t^2v_{t,0}^2=o_p(1),
\]
and the displayed variance convergence proves the lemma.
\end{proof}

\begin{lemma}[Oracle maximum bound]
\label{lem:wb-oracle-max}
Under Assumptions \ref{ass:prim-stoch}, \ref{ass:prim-bounded},
\ref{ass:prim-gram}, and \ref{ass:sol-oracle-envelope},
\[
    \max_{t\leq T}|U_tv_{t,0}|=o_p(\sqrt T).
\]
\end{lemma}

\begin{proof}[Detailed solution]
Conditional on \(\D_T\), sub-Gaussian maximal inequalities give
\[
    \max_{t\leq T}|U_t|=O_p(\sqrt{\log T}).
\]
Lemma \ref{lem:wb-leverage} and Assumption \ref{ass:sol-oracle-envelope} give
\(\max_t|v_{t,0}|=O_p(b_K)\).  Therefore
\[
    \max_{t\leq T}|U_tv_{t,0}|=O_p(b_K\sqrt{\log T})=o_p(\sqrt T),
\]
by \(b_K^2\log T/T\to0\).
\end{proof}

\begin{lemma}[Feasible true-partition replacement]
\label{lem:wb-feasible-replacement}
Under Assumptions \ref{ass:prim-stoch}--\ref{ass:prim-sieve} and
\ref{ass:sol-oracle-envelope}, suppose
\[
    T^{-1}\sum_{t=1}^T \sigma_t^2v_{t,0}^2\xrightarrow{p}\eta_q^2.
\]
Then
\begin{align}
    T^{-1/2}\sum_{t=1}^T\widehat Z_t(\beta_0,\Lambda_0)
    &=
    T^{-1/2}\sum_{t=1}^TU_tv_{t,0}+o_p(1),
    \\
    T^{-1}\sum_{t=1}^T\widehat Z_t(\beta_0,\Lambda_0)^2
    &=
    T^{-1}\sum_{t=1}^TU_t^2v_{t,0}^2+o_p(1),
    \\
    \max_{t\leq T}|\widehat Z_t(\beta_0,\Lambda_0)|
    &=o_p(\sqrt T).
\end{align}
\end{lemma}

\begin{proof}[Detailed solution]
Let \(\bm Q_0=\bm Q_{K,\Lambda_0}\),
\(\bm\Pi_0=\bm I_T-\bm M_{K,\Lambda_0}\), and
\(\bm v_0=\bm M_{K,\Lambda_0}\bm w\).  The sieve representation gives
\[
    \bm Y-\beta_0\bm X_{lag}
    =\bm Q_0\bm c_0+\bm r_0+\bm U.
\]
By Lemma \ref{lem:wb-projection},
\[
    \widehat{\bm u}(\beta_0,\Lambda_0)
    =\bm M_{K,\Lambda_0}(\bm r_0+\bm U).
\]
Define
\[
    \bm e_0=\bm M_{K,\Lambda_0}\bm r_0,
    \qquad
    \bm h_0=\bm\Pi_0\bm U,
    \qquad
    \bm d_0=(\bm e_0-\bm h_0)\odot\bm v_0 .
\]
Then
\[
    \widehat Z_t(\beta_0,\Lambda_0)=U_tv_{t,0}+d_{t,0}.
\]
Conditional on \(\D_T\), \(\bm Q_0'\bm U\) is a sub-Gaussian vector with
second moment of order \(Td_K\).  Hence
\[
    \|\bm Q_0'\bm U\|=O_p(\sqrt{Td_K}).
\]
The Gram bounds imply
\[
    \|\bm h_0\|_T=O_p(\sqrt{d_K/T}),
    \qquad
    \|\bm h_0\|_\infty=O_p\left(\zeta_K\sqrt{d_K/T}\right).
\]
Projection contraction and Assumption \ref{ass:prim-sieve} give
\[
    \|\bm e_0\|_T=O_p(a_K),
    \qquad
    \|\bm e_0\|_\infty=O_p\{a_K(1+\zeta_K)\}.
\]
Since \(\|\bm v_0\|_\infty=O_p(b_K)\),
\[
    \|\bm d_0\|_T
    =O_p\left(a_Kb_K+b_K\sqrt{d_K/T}\right)=o_p(1),
\]
and
\[
    \|\bm d_0\|_\infty
    =O_p\left(a_Kb_K^2+\zeta_Kb_K\sqrt{d_K/T}\right)=o_p(\sqrt T).
\]
The numerator has the exact cancellation
\[
    T^{-1/2}\sum_t\widehat Z_t(\beta_0,\Lambda_0)
    =T^{-1/2}\bm U'\bm v_0
      +T^{-1/2}\bm r_0'\bm M_{K,\Lambda_0}\bm w
    =T^{-1/2}\bm U'\bm v_0+o_p(1).
\]
For the denominator,
\[
    \left|
      T^{-1}\sum_t\widehat Z_t^2
      -T^{-1}\sum_tU_t^2v_{t,0}^2
    \right|
    \leq
      2\|\bm U\odot\bm v_0\|_T\|\bm d_0\|_T
      +\|\bm d_0\|_T^2=o_p(1),
\]
because Lemma \ref{lem:wb-oracle-lln} gives
\(\|\bm U\odot\bm v_0\|_T=O_p(1)\).  The maximum bound follows from Lemma
\ref{lem:wb-oracle-max} and \(\|\bm d_0\|_\infty=o_p(\sqrt T)\).
\end{proof}

\begin{theorem}[Known-partition conditional-design Wilks]
\label{thm:wb-known-wilks}
Under Assumptions \ref{ass:prim-stoch}--\ref{ass:prim-sieve},
\ref{ass:prim-convex}, and \ref{ass:sol-oracle-envelope}, suppose
\[
    T^{-1}\sum_{t=1}^T \sigma_t^2v_{t,0}^2\xrightarrow{p}\eta_q^2,
    \qquad
    \Pp(\eta_q\geq\sqrt{c_\eta})=1
\]
for some \(c_\eta>0\).  Then
\[
    \ell_T(\beta_0,\Lambda_0)\Rightarrow\chi_1^2.
\]
\end{theorem}

\begin{proof}[Detailed solution]
Let
\[
    A_T=T^{-1/2}\sum_{t=1}^T
       \widehat Z_t(\beta_0,\Lambda_0),
    \qquad
    B_T=T^{-1}\sum_{t=1}^T
       \widehat Z_t(\beta_0,\Lambda_0)^2.
\]
Lemma \ref{lem:wb-feasible-replacement}, followed by Lemmas
\ref{lem:wb-oracle-clt} and \ref{lem:wb-oracle-lln}, gives
\[
    A_T\xrightarrow{\mathrm{st}}\eta_qZ,
    \qquad
    B_T\xrightarrow{p}\eta_q^2.
\]
Because \(\eta_q^2\geq c_\eta\) almost surely, \(B_T\geq c_\eta/2\) with
probability tending to one.  Lemmas \ref{lem:wb-oracle-max} and
\ref{lem:wb-feasible-replacement} give
\[
    \max_t|\widehat Z_t(\beta_0,\Lambda_0)|=o_p(\sqrt T).
\]
Assumption \ref{ass:prim-convex} supplies the convex-hull condition.  Lemma
\ref{lem:sol-scalar-el} gives
\[
    \ell_T(\beta_0,\Lambda_0)=A_T^2/B_T+o_p(1),
\]
and stable Slutsky yields \(A_T^2/B_T\Rightarrow Z^2\sim\chi_1^2\).
\end{proof}

\begin{remark}[Proof audit for Theorem \ref{thm:wb-known-wilks}]
The known-partition result now uses a single stochastic route.  Conditional on
\(\D_T\), the residualized score weights are fixed, so the numerator CLT,
denominator LLN, and maximum bound are ordinary conditional independent-array
arguments.  Break-date localization and order selection are irrelevant until
\(\Lambda_0\) is replaced by an estimated partition.
\end{remark}
\section{One-Break Estimated-Partition Bridge}
\label{sec:one-break-bridge}

In this section only, take \(q_0=1\), \(\Lambda_0=(k_0)\), and \(\widetilde\Lambda=(\widetilde k)\). Let
\[
    \Delta_M=\bm M_{K,\widetilde\Lambda}-\bm M_{K,\Lambda_0},
    \qquad
    \mathcal M_T=\{t:j_{\widetilde\Lambda}(t)\neq j_{\Lambda_0}(t)\}.
\]

\begin{lemma}[Boundary-set size]
\label{lem:wb-boundary-size}
If \(|\widetilde k-k_0|=O_p(r_T)\), then
\[
    |\mathcal M_T|=O_p(r_T).
\]
\end{lemma}


\begin{proof}[Detailed solution]
Suppose first that \(\widetilde k\geq k_0\).  For \(t\leq k_0\), both
partitions assign observation \(t\) to the left regime; for
\(t>\widetilde k\), both assign it to the right regime.  The assignments
differ exactly on
\[
    k_0<t\leq\widetilde k.
\]
Thus \(|\mathcal M_T|=\widetilde k-k_0\).  If \(\widetilde k<k_0\), the
same argument gives the differing set \(\widetilde k<t\leq k_0\), so
\(|\mathcal M_T|=k_0-\widetilde k\).  Hence, in all cases,
\[
    |\mathcal M_T|=|\widetilde k-k_0|=O_p(r_T).
\]
\end{proof}
\begin{lemma}[Search-uniform projection noise]
\label{lem:wb-uniform-projection-noise}
Suppose \(q_0=1\) and Assumptions \ref{ass:prim-stoch},
\ref{ass:prim-bounded}, \ref{ass:prim-gram}, and \ref{ass:sol-local-series}
hold.  For every fixed \(C<\infty\), over the local one-break set
\[
    \mathcal N_T(Cr_T)=\{k:|k-k_0|\leq Cr_T\},
\]
\[
    \sup_{k\in\mathcal N_T(Cr_T)}\|\bm\Pi_k\bm U\|_T
    =O_p\left(\sqrt{\frac{K+\log(1+r_T)}{T}}\right),
\]
and
\[
    \sup_{k\in\mathcal N_T(Cr_T)}\|\bm\Pi_k\bm U\|_\infty
    =O_p\left(
      \zeta_K\sqrt{\frac{\log\{T(1+r_T)\}}{T}}
    \right).
\]
\end{lemma}

\begin{proof}[Detailed solution]
Condition on \(\D_T\).  For fixed \(k\), \(\bm\Pi_k\bm U\) is the projection
of a conditionally sub-Gaussian vector onto a subspace of rank at most \(2K\).
Standard chi-square/sub-Gaussian projection concentration gives
\[
    \|\bm\Pi_k\bm U\|^2=O_p(K+x)
\]
with conditional tail probability bounded by \(e^{-cx}\).  A union bound over
\(O(1+r_T)\) local candidates yields the displayed \(T\)-norm bound.  For the
sup norm, the leverage bound gives \((\bm\Pi_k)_{tt}\leq C\zeta_K^2/T\)
uniformly over \(t\) and local \(k\).  Conditional sub-Gaussian tails for each
coordinate, followed by a union bound over \(T(1+r_T)\) pairs \((t,k)\), give
the displayed sup-norm bound.
\end{proof}


\begin{lemma}[Nuisance part of partition perturbation]
\label{lem:wb-bridge-mu}
Suppose \(q_0=1\), Assumptions \ref{ass:prim-bounded},
\ref{ass:prim-gram}, \ref{ass:prim-sieve}, and
\ref{ass:sol-local-bias} hold, and \(|\widetilde k-k_0|=O_p(r_T)\).  Then
\[
    T^{-1/2}\bm\mu'\Delta_M\bm w
    =
    O_p\!\left(\frac{b_Kr_T}{\sqrt T}\right)+o_p(1),
    \qquad b_K=1+\zeta_K.
\]
In particular, this term is \(o_p(1)\) if \(b_Kr_T/\sqrt T\to0\).
\end{lemma}


\begin{proof}[Detailed solution]
It suffices to prove a uniform deterministic bound for all
\(|k-k_0|\leq Cr_T\), then evaluate it at \(k=\widetilde k\) on the
localization event.  Consider first \(k=k_0+h\), \(h\geq0\), and let
\(H=\{k_0+1,\ldots,k_0+h\}\).  Let \(\bm c_0\) and \(\bm c_1\) be
the true-regime sieve coefficients.  If the candidate partition uses
\(\bm c_0\) on its left block and \(\bm c_1\) on its right block, then
\[
    \bm\mu=\bm Q_{K,k}\bm c^{(k)}+\bm d_H+\bm r_0,
\]
where \(\bm d_H\) is supported on \(H\).  For \(t\in H\),
\[
    d_{H,t}
    =\bm P_K(W_{t-1})'(\bm c_1-\bm c_0)
    =m_1(W_{t-1})-m_0(W_{t-1})-
      \{r_{1,K,t}-r_{0,K,t}\}.
\]
Assumptions \ref{ass:prim-bounded} and \ref{ass:prim-sieve} imply
\(\|\bm d_H\|_\infty=O_p(1)\).  Since
\(\bm Q_{K,k}'\bm M_{K,k}=\bm0\),
\[
    \bm\mu'\bm M_{K,k}\bm w
    =\bm d_H'\bm M_{K,k}\bm w
      +\bm r_0'\bm M_{K,k}\bm w.
\]
At the true partition,
\[
    \bm\mu'\bm M_{K,k_0}\bm w
    =\bm r_0'\bm M_{K,k_0}\bm w.
\]
Therefore, using Lemma \ref{lem:wb-leverage} uniformly over local candidates,
\begin{align*}
    T^{-1/2}|\bm\mu'(\bm M_{K,k}-\bm M_{K,k_0})\bm w|
    &\leq
      T^{-1/2}\|\bm d_H\|_1\|\bm M_{K,k}\bm w\|_\infty \\
    &\quad+
      T^{-1/2}|\bm r_0'(\bm M_{K,k}-\bm M_{K,k_0})\bm w| \\
    &=O_p\left(\frac{b_Kh}{\sqrt T}\right)+o_p(1),
\end{align*}
where the final term is Assumption \ref{ass:sol-local-bias}.  The case
\(k<k_0\) is identical after exchanging the left and right regimes.  Since
\(|\widetilde k-k_0|=O_p(r_T)\), the lemma follows.

The sufficient primitive rate in \eqref{eq:sol-local-bias-sufficient} follows
from the same local coefficient update.  Write \(\bm v_k=\bm M_{K,k}\bm w\).
On the boundary set \(H\), \(|v_{t,k}-v_{t,k_0}|=O_p(b_K)\).  Off the
boundary, the candidate and true fitted-weight coefficients differ by
\(O_p(h\zeta_Kb_K/T)\), and each row has norm at most \(\zeta_K\).  Hence
\[
    \|\bm v_k-\bm v_{k_0}\|_1
    =O_p\{h b_K(1+\zeta_K^2)\}.
\]
Consequently,
\[
    T^{-1/2}|\bm r_0'(\bm v_k-\bm v_{k_0})|
    \leq T^{-1/2}\|\bm r_0\|_\infty
       \|\bm v_k-\bm v_{k_0}\|_1
    =O_p\left(
       \frac{a_Kb_K(1+\zeta_K^2)h}{\sqrt T}
      \right),
\]
which is uniformly \(o_p(1)\) for \(h\leq Cr_T\) under
\eqref{eq:sol-local-bias-sufficient}.
\end{proof}


\begin{lemma}[Stochastic part of partition perturbation]
\label{lem:wb-bridge-u}
Suppose \(q_0=1\), Assumptions \ref{ass:prim-stoch},
\ref{ass:prim-bounded}, \ref{ass:prim-gram}, and
\ref{ass:sol-local-series} hold, \(|\widetilde k-k_0|=O_p(r_T)\), and
\(b_Kr_T/\sqrt T\to0\).  Then
\[
    T^{-1/2}\bm U'\Delta_M\bm w
    =
    O_p\!\left(\zeta_K\sqrt{\frac{r_T\{K+\log(1+r_T)\}}{T}}\right)+o_p(1).
\]
In particular, this term is \(o_p(1)\) if
\(\zeta_K\sqrt{r_T\{K+\log(1+r_T)\}/T}\to0\).
\end{lemma}


\begin{proof}[Detailed solution]
Work uniformly over deterministic local candidates \(|k-k_0|\leq Cr_T\),
then evaluate at \(k=\widetilde k\) on the localization event.  Take
\(k=k_0+h>k_0\) and define
\[
    A=\{1,\ldots,k_0\},\quad
    H=\{k_0+1,\ldots,k_0+h\},\quad
    B=\{k_0+h+1,\ldots,T\},\quad
    R=H\cup B.
\]
For a segment \(J\), put
\[
    \bm G_J=\bm P_J'\bm P_J,
    \qquad
    \widehat{\bm\gamma}_J=\bm G_J^{-1}\bm P_J'\bm w_J.
\]
Assumption \ref{ass:sol-local-series} gives
\(\|\bm G_J^{-1}\|=O_p(T^{-1})\) for all long local segments, and bounded
weights imply \(\|\widehat{\bm\gamma}_J\|=O_p(1)\).  The exact update
identity
\[
    \widehat{\bm\gamma}_{A\cup H}-\widehat{\bm\gamma}_A
    =\bm G_{A\cup H}^{-1}
      \bm P_H'(\bm w_H-\bm P_H\widehat{\bm\gamma}_A)
\]
therefore yields
\[
    \|\widehat{\bm\gamma}_{A\cup H}-\widehat{\bm\gamma}_A\|
    =O_p\left(\frac{h\zeta_Kb_K}{T}\right).
\]
The same calculation after removing \(H\) from the true right segment gives
\[
    \|\widehat{\bm\gamma}_{B}-\widehat{\bm\gamma}_{R}\|
    =O_p\left(\frac{h\zeta_Kb_K}{T}\right).
\]

The residualized weights are segmentwise
\[
    \bm v_{k_0}=(\bm w_A-\bm P_A\widehat{\bm\gamma}_A,
             \bm w_R-\bm P_R\widehat{\bm\gamma}_R),
\]
\[
    \bm v_k=(\bm w_{A\cup H}-\bm P_{A\cup H}\widehat{\bm\gamma}_{A\cup H},
             \bm w_B-\bm P_B\widehat{\bm\gamma}_B).
\]
Cancelling the raw \(\bm w_H\) terms gives the exact decomposition
\begin{align}
    \bm U'(\bm v_k-\bm v_{k_0})
    &= (\bm P_A'\bm U_A)'
       (\widehat{\bm\gamma}_A-
        \widehat{\bm\gamma}_{A\cup H})
       \nonumber\\
    &\quad +(\bm P_B'\bm U_B)'
       (\widehat{\bm\gamma}_R-
        \widehat{\bm\gamma}_{B})
       \nonumber\\
    &\quad +(\bm P_H'\bm U_H)'
       (\widehat{\bm\gamma}_R-
        \widehat{\bm\gamma}_{A\cup H}).
    \label{eq:sol-u-bridge-decomp}
\end{align}
Conditional on \(\D_T\), sub-Gaussian concentration and the Gram bounds give
\[
    \|\bm P_A'\bm U_A\|=O_p(\sqrt{TK}),
    \qquad
    \|\bm P_B'\bm U_B\|=O_p(\sqrt{TK}).
\]
For the local moving block, a union bound over \(h\leq Cr_T\) gives
\[
    \sup_{h\leq Cr_T}\|\bm P_H'\bm U_H\|
      =O_p\{\zeta_K\sqrt{r_T\log(1+r_T)}\}.
\]
The coefficient difference in the last line of
\eqref{eq:sol-u-bridge-decomp} is \(O_p(1)\).  Combining the displays,
uniformly for \(h\leq Cr_T\),
\begin{align*}
    T^{-1/2}|\bm U'(\bm v_k-\bm v_{k_0})|
    &=O_p\left(\frac{h\zeta_Kb_K\sqrt K}{T}\right)
      +O_p\left(\zeta_K\sqrt{\frac{r_T\log(1+r_T)}{T}}\right).
\end{align*}
The first order is negligible under \(b_Kr_T/\sqrt T\to0\), and the second
order is exactly the logarithmic local-search term.  Taking \(h=O_p(r_T)\)
proves the displayed result.
The case \(k<k_0\) is the same with left and right interchanged.
\end{proof}


\begin{lemma}[Denominator perturbation]
\label{lem:wb-bridge-v}
Suppose \(q_0=1\), Assumptions \ref{ass:prim-stoch},
\ref{ass:prim-score-admissible}, \ref{ass:prim-bounded},
\ref{ass:prim-gram}, \ref{ass:prim-sieve}, \ref{ass:sol-oracle-envelope},
\ref{ass:sol-local-bias}, and \ref{ass:sol-local-series} hold.  Suppose also
\[
    T^{-1}\sum_{t=1}^T\sigma_t^2v_{t,0}^2\xrightarrow{p}\eta_q^2,
    \qquad |\widetilde k-k_0|=O_p(r_T),
\]
and
\[
    \frac{b_Kr_T}{\sqrt T}\to0,
    \qquad
    b_K\sqrt{\frac{K+\log(1+r_T)}{T}}\to0,
    \qquad
    \frac{\zeta_K^2K r_T}{T}\to0.
\]
Then
\[
    V_T(\widetilde\Lambda)-V_T(\Lambda_0)
    =O_p\left(
      \frac{b_K^2r_T}{T}
      +\frac{\zeta_K^2K r_T}{T}
    \right)+o_p(1)=o_p(1).
\]
\end{lemma}


\begin{proof}[Detailed solution]
Again work uniformly over \(|k-k_0|\leq Cr_T\).  Put
\(\bm v_k=\bm M_{K,k}\bm w\) and
\(V_{U,T}(k)=T^{-1}\sum_tU_t^2v_{t,k}^2\).  As in Lemma
\ref{lem:wb-bridge-mu}, for a local candidate partition,
\[
    \bm\mu=\bm Q_{K,k}\bm c^{(k)}+\bm d_H+\bm r_0,
\]
where \(\bm d_H\) is supported on the boundary set and
\(\|\bm d_H\|_\infty=O_p(1)\).  Therefore
\[
    \widehat{\bm u}(\beta_0,k)
    =\bm U-\bm\Pi_k\bm U+\bm M_{K,k}\bm r_0+\bm M_{K,k}\bm d_H,
    \qquad \bm\Pi_k=\bm I_T-\bm M_{K,k},
\]
and
\[
    \widehat{\bm Z}(\beta_0,k)
    =\bm U\odot\bm v_k+\bm\rho_k,
\]
with
\[
    \|\bm\rho_k\|_T
    \leq \|\bm v_k\|_\infty
       \{\|\bm\Pi_k\bm U\|_T
        +\|\bm M_{K,k}\bm r_0\|_T
        +\|\bm M_{K,k}\bm d_H\|_T\}.
\]
Lemma \ref{lem:wb-uniform-projection-noise} gives, uniformly over the local
full-rank set,
\[
    \|\bm\Pi_k\bm U\|_T
    =O_p\left(\sqrt{\frac{K+\log(1+r_T)}{T}}\right).
\]
Projection contraction
gives
\[
    \|\bm M_{K,k}\bm r_0\|_T=O_p(a_K),
    \qquad
    \|\bm M_{K,k}\bm d_H\|_T=O_p(\sqrt{h/T}).
\]
Since \(\|\bm v_k\|_\infty=O_p(b_K)\), Assumptions
\ref{ass:prim-sieve} and \ref{ass:sol-oracle-envelope}, together with
\(b_Kr_T/\sqrt T\to0\), imply
\begin{equation}
    \sup_{|k-k_0|\leq Cr_T}\|\bm\rho_k\|_T=o_p(1).
    \label{eq:sol-rho-local}
\end{equation}

It remains to compare the oracle denominators.  The coefficient-update bounds
from Lemma \ref{lem:wb-bridge-u} imply, for \(h=|k-k_0|\),
\[
    |v_{t,k}-v_{t,k_0}|
    \leq
    \begin{cases}
      Cb_K, & t\in\mathcal M_T,\\[1mm]
      C h\zeta_K^2b_K/T, & t\notin\mathcal M_T.
    \end{cases}
\]
Also \(\max_t(|v_{t,k}|+|v_{t,k_0}|)=O_p(b_K)\).  For a boundary interval
of length \(h\), conditional sub-Gaussian concentration and a dyadic union
bound give
\[
    \sup_{1\leq h\leq Cr_T}\frac{1}{h}\sum_{t\in H_h}U_t^2=O_p(1),
\]
and \(T^{-1}\sum_tU_t^2=O_p(1)\).  Hence, uniformly over local candidates,
\begin{align*}
    |V_{U,T}(k)-V_{U,T}(k_0)|
    &\leq T^{-1}\sum_tU_t^2
      |v_{t,k}-v_{t,k_0}|(|v_{t,k}|+|v_{t,k_0}|)\\
    &=O_p\left(\frac{b_K^2h}{T}\right)
      +O_p\left(\frac{h\zeta_K^2b_K^2}{T}\right).
\end{align*}
Assumption \ref{ass:sol-local-series} gives \(b_K^2=O(K)\), so the second
term is \(O_p(\zeta_K^2Kh/T)\).  Thus
\[
    |V_{U,T}(k)-V_{U,T}(k_0)|
    =O_p\left(\frac{b_K^2h}{T}+\frac{\zeta_K^2Kh}{T}\right).
\]
By Lemma \ref{lem:wb-oracle-lln}, \(V_{U,T}(k_0)=O_p(1)\); hence
\(\|\bm U\odot\bm v_k\|_T=O_p(1)\) uniformly locally.  The
square-difference inequality and \eqref{eq:sol-rho-local} imply
\[
    V_T(k)-V_{U,T}(k)=o_p(1)
\]
uniformly locally, and the same is true at \(k_0\).  Taking
\(h=|\widetilde k-k_0|=O_p(r_T)\) gives the displayed bound.
\end{proof}


\begin{theorem}[One-break estimated-partition score equivalence]
\label{thm:wb-one-break-equivalence}
Suppose \(q_0=1\), Assumptions \ref{ass:prim-stoch}--\ref{ass:prim-localization-target},
\ref{ass:sol-oracle-envelope}, \ref{ass:sol-local-bias}, and
\ref{ass:sol-local-series} hold, and
\[
    T^{-1}\sum_{t=1}^T\sigma_t^2v_{t,0}^2\xrightarrow{p}\eta_q^2.
\]
Assume \(r_T\geq1\) and the bridge rates
\[
    \zeta_K\sqrt{\frac{r_T\{K+\log(1+r_T)\}}{T}}\to0,
    \qquad
    \frac{b_Kr_T}{\sqrt T}\to0,
    \qquad
    \frac{\zeta_K^2K r_T}{T}\to0.
\]
Then
\[
    S_T(\widetilde\Lambda)-S_T(\Lambda_0)=o_p(1),
    \qquad
    V_T(\widetilde\Lambda)-V_T(\Lambda_0)=o_p(1).
\]
\end{theorem}


\begin{proof}[Detailed solution]
Assumption \ref{ass:prim-localization-target} gives
\(|\widetilde k-k_0|=O_p(r_T)\).  For every one-break partition
\(\Lambda\), symmetry and idempotence imply
\[
    S_T(\Lambda)
    =T^{-1/2}(\bm U+\bm\mu)'\bm M_{K,\Lambda}\bm w.
\]
Therefore
\[
    S_T(\widetilde\Lambda)-S_T(\Lambda_0)
    =T^{-1/2}\bm\mu'\Delta_M\bm w
      +T^{-1/2}\bm U'\Delta_M\bm w.
\]
Lemma \ref{lem:wb-bridge-mu} bounds the nuisance term by
\[
    O_p\left(\frac{b_Kr_T}{\sqrt T}\right)+o_p(1)=o_p(1),
\]
and Lemma \ref{lem:wb-bridge-u} bounds the stochastic term by
\[
    O_p\left(\zeta_K\sqrt{\frac{r_T\{K+\log(1+r_T)\}}{T}}\right)+o_p(1)=o_p(1).
\]
This proves the numerator equivalence.

For the denominator, Lemma \ref{lem:wb-bridge-v} gives
\[
    V_T(\widetilde\Lambda)-V_T(\Lambda_0)
    =O_p\left(
       \frac{b_K^2r_T}{T}+
       \frac{\zeta_K^2Kr_T}{T}
      \right)+o_p(1).
\]
The third bridge rate makes the second term \(o_p(1)\).  Since \(r_T\geq1\),
\[
    \frac{b_K^2r_T}{T}
    =\frac{1}{r_T}
      \left(\frac{b_Kr_T}{\sqrt T}\right)^2=o(1),
\]
so the denominator difference is \(o_p(1)\).
\end{proof}

\begin{remark}[Proof audit for Theorem \ref{thm:wb-one-break-equivalence}]
Theorem \ref{thm:wb-one-break-equivalence} is only a bridge from the true
one-break partition to a localized estimated break.  Lemma
\ref{lem:wb-boundary-size} limits the number of observations whose regime
assignment can change.  Lemma \ref{lem:wb-bridge-mu} handles the deterministic
nuisance component and makes the local sieve-bias condition explicit.  Lemma
\ref{lem:wb-bridge-u} handles the stochastic component through conditional
sub-Gaussian concentration and segmentwise coefficient updates.  Lemma \ref{lem:wb-bridge-v} proves that the empirical
second moment is stable after the same partition perturbation.  No convex-hull
condition and no empirical-likelihood expansion is used in this theorem; those
enter only in the next Wilks theorem.  The result would fail without
localization, local full-rank Gram regularity, and the three bridge rates
listed in the theorem statement.
\end{remark}

\begin{theorem}[One-break primitive estimated-partition Wilks]
\label{thm:wb-one-break-wilks}
Under the assumptions of Theorem \ref{thm:wb-known-wilks}, the convex-hull condition for the estimated-partition score array, and Theorem \ref{thm:wb-one-break-equivalence},
\[
    \ell_T(\beta_0,\widetilde\Lambda)\Rightarrow\chi_1^2.
\]
\end{theorem}


\begin{proof}[Detailed solution]
Let \(A_T(\Lambda)=S_T(\Lambda)\) and \(B_T(\Lambda)=V_T(\Lambda)\).
Theorem \ref{thm:wb-one-break-equivalence} gives
\[
    A_T(\widetilde\Lambda)=A_T(\Lambda_0)+o_p(1),
    \qquad
    B_T(\widetilde\Lambda)=B_T(\Lambda_0)+o_p(1).
\]
The known-partition proof gives
\[
    A_T(\Lambda_0)\xrightarrow{\mathrm{st}}\eta_qZ,
    \qquad
    B_T(\Lambda_0)\xrightarrow{p}\eta_q^2.
\]
We must also verify the maximum condition; numerator and denominator
 equivalence alone would not imply it.  The local representation in the proof
of Lemma \ref{lem:wb-bridge-v} gives
\[
    \widehat{\bm Z}(\beta_0,\widetilde\Lambda)
    =\bm U\odot\bm v_{\widetilde k}+\bm\rho_{\widetilde k}.
\]
The envelope rate implies
\(\max_t|U_tv_{t,\widetilde k}|=o_p(\sqrt T)\).  For the remainder, the
uniform projection calculation gives
\[
    \|\bm\Pi_{\widetilde k}\bm U\|_\infty
    =O_p\left(
      \zeta_K\sqrt{\frac{\log\{T(1+r_T)\}}{T}}
    \right),
\]
while the sieve calculation gives
\[
    \|\bm M_{K,\widetilde k}\bm r_0\|_\infty
    =O_p\{a_K(1+\zeta_K)\}.
\]
Finally, because \(\bm d_H\) has bounded entries and support of size
\(O_p(r_T)\),
\[
    \|\bm M_{K,\widetilde k}\bm d_H\|_\infty
    =O_p\left(1+\frac{r_T\zeta_K^2}{T}\right)=O_p(1).
\]
Multiplication by \(\|\bm v_{\widetilde k}\|_\infty=O_p(b_K)\), followed by
the sieve, bridge, and envelope rates, gives
\(\|\bm\rho_{\widetilde k}\|_\infty=o_p(\sqrt T)\).  Hence
\[
    \max_t|\widehat Z_t(\beta_0,\widetilde\Lambda)|=o_p(\sqrt T).
\]
The estimated-partition convex-hull condition and Lemma
\ref{lem:sol-scalar-el} now yield
\[
    \ell_T(\beta_0,\widetilde\Lambda)
    =\frac{A_T(\widetilde\Lambda)^2}
           {B_T(\widetilde\Lambda)}+o_p(1).
\]
Substitution of the two equivalences reduces this ratio to the
true-partition ratio plus \(o_p(1)\), and Theorem
\ref{thm:wb-known-wilks} gives the \(\chi_1^2\) limit.
\end{proof}


\section{Conditional-Design Profile-Sieve Segmentation}
\label{sec:break-localization}

This section is the main rebuild of the Bai--Perron module.  The break search
is no longer imported as a broad high-level assumption.  Instead, the stochastic
part is controlled by conditional-on-design sub-Gaussian concentration and the
identifying part is isolated as a deterministic profile-loss condition.

At the null value \(\beta_0\), define
\[
    y_t(\beta_0)=Y_t-\beta_0X_{t-1}.
\]
Then the nuisance-break model and sieve approximation imply
\[
    y_t(\beta_0)
    =
    \bm P_K(W_{t-1})'\bm c_{j,K}+U_t+r_{j,K,t},
    \qquad
    k_{j,0}<t\leq k_{j+1,0}.
\]
Thus, apart from approximation error, the break search is a segmented least-
squares problem with regime-specific sieve coefficient vectors.

\begin{assumption}[Conditional-design profile-sieve segmentation]
\label{ass:bp-profile-sieve}
For the main estimated-partition theory, \(q_0=1\), \(\Lambda_0=(k_0)\), and
\(\bar q\) is fixed.  Let \(\bm\mu_K=\bm Q_{K,\Lambda_0}\bm c_0\) be the sieve
part of the true conditional mean and define
\[
    D_T(\Lambda)=\|\bm M_{K,\Lambda}\bm\mu_K\|^2.
\]
The following conditions hold.
\begin{enumerate}
    \item \textbf{Admissible partitions.}  All searches are over
    \(\mathcal L_q(\epsilon)\), the true break obeys the stricter spacing
    \(\epsilon_0>\epsilon\), and \((\bar q+1)\epsilon<1\).

    \item \textbf{Global segment conditioning.}  The global Gram condition in
    Assumption \ref{ass:prim-gram} holds for every admissible segment used in
    the search.

    \item \textbf{Deterministic profile identification.}  For every admissible
    one-break candidate,
    \[
        D_T(k)\geq c_\Delta |k-k_0|\Delta_T^2.
    \]
    For the underfitted zero-break model,
    \[
        \inf_{\Lambda\in\mathcal L_0(\epsilon)}D_T(\Lambda)
        \geq c_\Delta T\Delta_T^2.
    \]

    \item \textbf{Uniform conditional RSS expansion.}  With
    \[
        R_T=K+\log T+Ta_K^2,
        \qquad
        r_T=R_T/\Delta_T^2,
    \]
    conditional sub-Gaussian concentration gives, uniformly over the admissible
    one-break search,
    \[
        RSS_K(\beta_0,k)-RSS_K(\beta_0,k_0)
        \geq \frac12D_T(k)-O_p(R_T),
    \]
    and uniformly over \(0\leq q\leq\bar q\),
    \[
        RSS_K(\beta_0,k_0)
        -RSS_K\{\beta_0,\widehat\Lambda_q(\beta_0)\}=O_p(R_T)
        \quad(q>1).
    \]
    The rates satisfy
    \[
        R_T=o(T\Delta_T^2),
        \qquad
        r_T=o(\sqrt T).
    \]
\end{enumerate}
\end{assumption}

\begin{proposition}[Profile-sieve reduction to segmented least squares]
\label{prop:bp-profile-reduction}
At \(\beta_0\), the fixed-order nuisance-break estimator is the least-squares
segmentation estimator in the generated regression
\[
    y_t(\beta_0)=\bm P_K(W_{t-1})'\bm c_j+\text{error},
    \qquad t\in I_j(\Lambda).
\]
\end{proposition}

\begin{proof}[Detailed solution]
For a fixed candidate partition \(\Lambda\), minimizing the null-imposed
segment cost over the nuisance coefficients gives
\[
    RSS_K(\beta_0,\Lambda)
    =
    \sum_{j=0}^{q}
    \min_{\bm c_j\in\R^K}
    \sum_{t=k_j+1}^{k_{j+1}}
    \{y_t(\beta_0)-\bm P_K(W_{t-1})'\bm c_j\}^2.
\]
This is the segmented least-squares criterion for a linear regression whose
coefficient vector may change across regimes.  The block matrix
\(\bm Q_{K,\Lambda}\) is the stacked design matrix and
\(\bm M_{K,\Lambda}\) is its residual-maker.  The nonstandard work is exactly
what Assumption \ref{ass:bp-profile-sieve} records: growing \(K\), approximation
error, uniform segment conditioning, and conditional concentration over the
finite admissible search.
\end{proof}

\begin{lemma}[One-break profile-RSS exclusion]
\label{lem:wb-bp-fixed-exclusion}
Under Assumption \ref{ass:bp-profile-sieve}, for every fixed \(M\), define
\[
    \mathcal B_T(M)
    =\{k\in\mathcal L_1(\epsilon): |k-k_0|\geq Mr_T\}.
\]
Then
\[
    \lim_{M\to\infty}\liminf_{T\to\infty}
    \Pp\left\{
      \inf_{k\in\mathcal B_T(M)}
      [RSS_K(\beta_0,k)-RSS_K(\beta_0,k_0)]>0
    \right\}=1.
\]
\end{lemma}

\begin{proof}[Detailed solution]
For \(k\in\mathcal B_T(M)\), deterministic identification and the uniform RSS
expansion give
\[
    RSS_K(\beta_0,k)-RSS_K(\beta_0,k_0)
    \geq \frac12 c_\Delta M r_T\Delta_T^2-O_p(R_T)
    =\left(\frac12c_\Delta M+O_p(1)\right)R_T.
\]
Taking \(T\to\infty\) and then \(M\to\infty\) proves the exclusion event.
\end{proof}

\begin{theorem}[One-break profile-sieve localization]
\label{thm:wb-localization}
Under Assumption \ref{ass:bp-profile-sieve},
\[
    |\widehat k_1(\beta_0)-k_0|=O_p(r_T),
    \qquad
    r_T=R_T/\Delta_T^2=o(\sqrt T).
\]
\end{theorem}

\begin{proof}[Detailed solution]
The one-break estimator minimizes \(RSS_K(\beta_0,k)\) over
\(\mathcal L_1(\epsilon)\), and \(k_0\) is admissible.  On the exclusion event
in Lemma \ref{lem:wb-bp-fixed-exclusion}, no candidate with
\(|k-k_0|\geq Mr_T\) can minimize the criterion.  Hence
\[
    \Pp\{|\widehat k_1(\beta_0)-k_0|\geq Mr_T\}
    \leq
    \Pp\left\{
      \inf_{k\in\mathcal B_T(M)}
      [RSS_K(\beta_0,k)-RSS_K(\beta_0,k_0)]\leq0
    \right\},
\]
and the right side vanishes after taking \(T\to\infty\) and then
\(M\to\infty\).
\end{proof}

\begin{assumption}[Concrete unknown-order penalty]
\label{ass:wb-penalty}
Use
\[
    \operatorname{pen}_T(q,K)=q\kappa_T R_T,
\]
where
\[
    \kappa_T\to\infty,
    \qquad
    \kappa_T R_T=o(T\Delta_T^2).
\]
For instance, \(\kappa_T=\log\log(T+e^e)\) is valid whenever
\(R_T\log\log T=o(T\Delta_T^2)\).
\end{assumption}

\begin{theorem}[Unknown-order consistency, one-break truth]
\label{thm:wb-order}
Under Assumptions \ref{ass:bp-profile-sieve} and \ref{ass:wb-penalty}, with
\(q_0=1\),
\[
    \Pp\{\widehat q(\beta_0)=1\}\to1.
\]
\end{theorem}

\begin{proof}[Detailed solution]
Write \(RSS_q^*=RSS_K\{\beta_0,\widehat\Lambda_q(\beta_0)\}\).  For the
underfitted order \(q=0\), Assumption \ref{ass:bp-profile-sieve} gives
\[
    RSS_0^*-RSS_1^*
    \geq cT\Delta_T^2-O_p(R_T).
\]
The penalty advantage of underfitting is only \(\kappa_TR_T=o(T\Delta_T^2)\),
so the penalized zero-break objective is larger with probability tending to one.
For every \(q>1\), the overfit RSS gain is \(O_p(R_T)\), while the penalty
increase is \((q-1)\kappa_TR_T\).  Since \(\kappa_T\to\infty\) and \(\bar q\)
is fixed, all overfitted objectives are larger with probability tending to one.
A finite union bound over \(q\neq1\) proves the claim.
\end{proof}
\begin{lemma}[High-probability event transfer]
\label{lem:wb-event-transfer}
Let \(X_T\) and \(Y_T\) be real statistics, and let \(E_T\) be events
such that \(\Pp(E_T)\to1\).  If \(X_T\Rightarrow X\) and
\(Y_T=X_T\) on \(E_T\), then \(Y_T\Rightarrow X\).
\end{lemma}

\begin{proof}[Detailed solution]
For every bounded continuous function \(f\),
\[
    |\E f(Y_T)-\E f(X_T)|
    \leq 2\|f\|_\infty\Pp(E_T^c)\to0.
\]
Since \(\E f(X_T)\to\E f(X)\), the same convergence holds for
\(\E f(Y_T)\), which is weak convergence of \(Y_T\) to \(X\).
\end{proof}

\begin{corollary}[Unknown-order Wilks, one-break truth]
\label{cor:wb-unknown-wilks}
Suppose \(q_0=1\), Theorem \ref{thm:wb-order} holds, the fixed-order
estimated-partition Wilks theorem \ref{thm:wb-one-break-wilks} holds for
\(\widetilde\Lambda=\widehat\Lambda_{1}(\beta_0)\), and the same measurable
tie-breaking convention is used for fixed-order and unknown-order selection.
Then
\[
    \ell_T\{\beta_0,\widehat\Lambda(\beta_0)\}\Rightarrow\chi_1^2.
\]
\end{corollary}

\begin{proof}[Detailed solution]
Let \(E_T=\{\widehat q(\beta_0)=1\}\).  By Theorem
\ref{thm:wb-order}, \(\Pp(E_T)\to1\).  On \(E_T\), unknown-order
selection uses the one-break order and the same deterministic tie-breaking
rule as the fixed-order estimator, so
\[
    \widehat\Lambda(\beta_0)=\widehat\Lambda_{1}(\beta_0)
    =\widetilde\Lambda.
\]
Therefore
\[
    \ell_T\{\beta_0,\widehat\Lambda(\beta_0)\}
    =\ell_T(\beta_0,\widetilde\Lambda)
\]
on \(E_T\).  The fixed-order estimated-partition Wilks theorem gives
\( \ell_T(\beta_0,\widetilde\Lambda)\Rightarrow\chi_1^2 \), and Lemma
\ref{lem:wb-event-transfer} transfers this limit to the unknown-order
statistic.
\end{proof}

\begin{remark}[Proof audit for Corollary \ref{cor:wb-unknown-wilks}]
The corollary uses only order consistency and event transfer, and it is stated
only for the one-break truth covered by Theorem \ref{thm:wb-one-break-wilks}.
It does not introduce a new empirical-likelihood expansion.  The deterministic
tie-breaking convention is essential: without it, the selected partition on
\(E_T\) need not be literally the same measurable object as
\(\widehat\Lambda_1(\beta_0)\).
\end{remark}
\section{Local Alternatives}
\label{sec:local-alternatives}

\begin{assumption}[Canonical local alternative and drift]
\label{ass:wb-local-oracle}
Let
\[
    \beta_T=\beta_0+\frac{h}{\sqrt T},
    \qquad h\in\R
\]
with deterministic \(h\).  Define
\[
    \Gamma_T
    =T^{-1}\bm X_{lag}'\bm M_{K,\Lambda_0}\bm w
    =T^{-1}\sum_{t=1}^T
      (\bm M_{K,\Lambda_0}\bm X_{lag})_t v_{t,0}.
\]
Under \(P_{\beta_T}\), Assumption \ref{ass:prim-stoch} and the Gram/envelope
conditions continue to hold, and
\[
    \Gamma_T\xrightarrow{p}\Gamma,
    \qquad
    \eta_T^2:=T^{-1}\sum_{t=1}^T\sigma_t^2v_{t,0}^2
    \xrightarrow{p}\eta^2,
    \qquad
    0<\eta^2<\infty,
\]
where \(\Gamma\) and \(\eta\) are deterministic.  Also
\[
    T^{-1}\sum_{t=1}^T
    \{(\bm M_{K,\Lambda_0}\bm X_{lag})_t v_{t,0}\}^2=O_p(1),
    \qquad
    \max_t |(\bm M_{K,\Lambda_0}\bm X_{lag})_t v_{t,0}|=o_p(T).
\]
Put
\[
    g_{t,T}=\frac{h}{\sqrt T}
    (\bm M_{K,\Lambda_0}\bm X_{lag})_t v_{t,0},
    \qquad
    \delta_h=\frac{h\Gamma}{\eta}.
\]
\end{assumption}

\begin{lemma}[Local shifted-score replacement]
\label{lem:wb-local-replacement}
Suppose Assumptions \ref{ass:prim-stoch}, \ref{ass:prim-sieve},
\ref{ass:prim-bounded}, \ref{ass:prim-gram}, and
\ref{ass:sol-oracle-envelope} continue to hold under \(P_{\beta_T}\).  Then
\begin{align}
    T^{-1/2}\sum_{t=1}^T
    \{\widehat Z_t(\beta_0,\Lambda_0)-(U_tv_{t,0}+g_{t,T})\}
    &=o_p(1),
    \\
    T^{-1}\sum_{t=1}^T
    \{\widehat Z_t(\beta_0,\Lambda_0)^2-(U_tv_{t,0}+g_{t,T})^2\}
    &=o_p(1),
    \\
    \max_t|\widehat Z_t(\beta_0,\Lambda_0)-(U_tv_{t,0}+g_{t,T})|
    &=o_p(\sqrt T).
\end{align}
\end{lemma}


\begin{proof}[Detailed solution]
Under \(P_{\beta_T}\),
\[
    \bm Y-\beta_0\bm X_{lag}
    =\bm Q_{K,\Lambda_0}\bm c_0+
      \bm r_0+\bm U+\frac{h}{\sqrt T}\bm X_{lag}.
\]
Apply \(\bm M_{K,\Lambda_0}\) and put
\[
    \bm e_0=\bm M_{K,\Lambda_0}\bm r_0,
    \qquad
    \bm h_0=\bm\Pi_{K,\Lambda_0}\bm U,
    \qquad
    \bm x_0^c=\bm M_{K,\Lambda_0}\bm X_{lag}.
\]
Then
\[
    \widehat{\bm u}(\beta_0,\Lambda_0)
    =\bm U-\bm h_0+\bm e_0+\frac{h}{\sqrt T}\bm x_0^c.
\]
Multiplication by \(\bm v_0=\bm M_{K,\Lambda_0}\bm w\) gives the exact
identity
\[
    \widehat Z_t(\beta_0,\Lambda_0)
    =U_tv_{t,0}+g_{t,T}+e^{\mathrm{feas}}_{t,T},
    \qquad
    e^{\mathrm{feas}}_{t,T}=(e_{t,0}-h_{t,0})v_{t,0}.
\]
The vector \(\bm e_T^{\mathrm{feas}}\) is the same feasibility error controlled
in Lemma \ref{lem:wb-feasible-replacement}; under the stated local-law
continuation of the primitive assumptions,
\[
    T^{-1/2}\sum_te^{\mathrm{feas}}_{t,T}=o_p(1),
    \qquad
    \|\bm e_T^{\mathrm{feas}}\|_T=o_p(1),
    \qquad
    \|\bm e_T^{\mathrm{feas}}\|_\infty=o_p(\sqrt T).
\]
The first and third displays of the lemma follow immediately.

For the square replacement, set \(a_{t,T}=U_tv_{t,0}+g_{t,T}\).  Then
\[
    T^{-1}\sum_t\{\widehat Z_t^2-a_{t,T}^2\}
    =2T^{-1}\sum_ta_{t,T}e^{\mathrm{feas}}_{t,T}
      +T^{-1}\sum_t(e^{\mathrm{feas}}_{t,T})^2.
\]
Assumption \ref{ass:wb-local-oracle} gives
\(\|\bm U\odot\bm v_0\|_T=O_p(1)\) and
\(\|\bm g_T\|_T=o_p(1)\), hence \(\|\bm a_T\|_T=O_p(1)\).  By
Cauchy--Schwarz and \(\|\bm e_T^{\mathrm{feas}}\|_T=o_p(1)\), both terms on the right are
\(o_p(1)\).
\end{proof}


\begin{lemma}[True-partition local empirical-likelihood limit]
\label{lem:wb-local-true-el}
Under Assumption \ref{ass:wb-local-oracle}, Lemma
\ref{lem:wb-local-replacement}, and the local convex-hull condition for
\(\widehat Z_t(\beta_0,\Lambda_0)\),
\[
    \ell_T(\beta_0,\Lambda_0)
    \Rightarrow\chi_1^2(\delta_h^2).
\]
\end{lemma}

\begin{proof}[Detailed solution]
Lemma \ref{lem:wb-local-replacement}, the conditional-design CLT from Lemma
\ref{lem:wb-oracle-clt} under the local law, and Assumption
\ref{ass:wb-local-oracle} give
\begin{align*}
    T^{-1/2}\sum_t\widehat Z_t(\beta_0,\Lambda_0)
    &=T^{-1/2}\sum_tU_tv_{t,0}
      +T^{-1/2}\sum_tg_{t,T}+o_p(1)\\
    &\xrightarrow{\mathrm{st}}\eta Z+h\Gamma
      =\eta(Z+\delta_h).
\end{align*}
For the second moment,
\begin{align*}
    T^{-1}\sum_t(U_tv_{t,0}+g_{t,T})^2
    &=T^{-1}\sum_tU_t^2v_{t,0}^2
      +2T^{-1}\sum_tU_tv_{t,0}g_{t,T}
      +T^{-1}\sum_tg_{t,T}^2.
\end{align*}
The first term converges to \(\eta^2\) and the last term is \(o_p(1)\).
The cross term is \(o_p(1)\) by Cauchy--Schwarz:
\[
    \left|T^{-1}\sum_tU_tv_{t,0}g_{t,T}\right|
    \leq
    \left(T^{-1}\sum_tU_t^2v_{t,0}^2\right)^{1/2}
    \left(T^{-1}\sum_tg_{t,T}^2\right)^{1/2}=o_p(1).
\]
The local replacement lemma therefore gives
\[
    T^{-1}\sum_t\widehat Z_t(\beta_0,\Lambda_0)^2
    \xrightarrow{p}\eta^2.
\]
The maximum condition follows from
\[
    \max_t|U_tv_{t,0}+g_{t,T}|=o_p(\sqrt T)
\]
and the maximum replacement bound in Lemma \ref{lem:wb-local-replacement}.
The scalar empirical-likelihood expansion, Lemma \ref{lem:sol-scalar-el}, now
yields
\[
    \ell_T(\beta_0,\Lambda_0)
    =\frac{\{\eta(Z+\delta_h)\}^2}{\eta^2}+o_p(1)
    \Rightarrow (Z+\delta_h)^2\sim\chi_1^2(\delta_h^2).
\]
\end{proof}


\begin{theorem}[Local noncentral chi-square]
\label{thm:wb-local-power}
Suppose the conditions of Lemma \ref{lem:wb-local-true-el} hold.  Suppose
also that, under \(P_{\beta_T}\), the local estimated-partition bridge holds:
\[
    S_T(\widetilde\Lambda)-S_T(\Lambda_0)=o_p(1),
    \qquad
    V_T(\widetilde\Lambda)-V_T(\Lambda_0)=o_p(1),
    \qquad
    \max_t|\widehat Z_t(\beta_0,\widetilde\Lambda)|=o_p(\sqrt T),
\]
and the local convex-hull condition holds for
\(\widehat Z_t(\beta_0,\widetilde\Lambda)\).  Then
\[
    \ell_T(\beta_0,\widetilde\Lambda)
    \Rightarrow
    \chi_1^2(\delta_h^2).
\]
If \(q_0=1\) and \(P_{\beta_T}\{\widehat q(\beta_0)=1\}\to1\), the same
limit holds for \(\ell_T\{\beta_0,\widehat\Lambda(\beta_0)\}\).
\end{theorem}


\begin{proof}[Detailed solution]
The proof of Lemma \ref{lem:wb-local-true-el} gives the true-partition score
limits
\[
    S_T(\Lambda_0)\xrightarrow{\mathrm{st}}\eta(Z+\delta_h),
    \qquad
    V_T(\Lambda_0)\xrightarrow{p}\eta^2,
    \qquad
    \max_t|\widehat Z_t(\beta_0,\Lambda_0)|=o_p(\sqrt T).
\]
The local bridge transfers the first two limits and supplies the maximum bound
for the estimated partition:
\[
    S_T(\widetilde\Lambda)\xrightarrow{\mathrm{st}}\eta(Z+\delta_h),
    \qquad
    V_T(\widetilde\Lambda)\xrightarrow{p}\eta^2,
    \qquad
    \max_t|\widehat Z_t(\beta_0,\widetilde\Lambda)|=o_p(\sqrt T).
\]
The local convex-hull condition and Lemma \ref{lem:sol-scalar-el} give
\[
    \ell_T(\beta_0,\widetilde\Lambda)
    =\frac{S_T(\widetilde\Lambda)^2}{V_T(\widetilde\Lambda)}+o_p(1)
    \Rightarrow (Z+\delta_h)^2\sim\chi_1^2(\delta_h^2).
\]
For unknown order, let \(E_T=\{\widehat q(\beta_0)=1\}\).  On \(E_T\),
\(\widehat\Lambda(\beta_0)=\widetilde\Lambda\) under the same tie-breaking
convention used in Corollary \ref{cor:wb-unknown-wilks}.  Since
\(P_{\beta_T}(E_T)\to1\), the event-transfer lemma
\ref{lem:wb-event-transfer} transfers the fixed-order local limit to the
unknown-order statistic.
\end{proof}

\begin{remark}[Proof audit for Theorem \ref{thm:wb-local-power}]
The local theorem has two separate pieces.  Lemma \ref{lem:wb-local-true-el}
proves the noncentral limit at the true partition from the local drift
\(g_{t,T}\), the local denominator stability, and the scalar EL expansion.
The theorem then uses a local-law version of the estimated-partition bridge to
move from \(\Lambda_0\) to \(\widetilde\Lambda\).  The bridge must include
the maximum-score condition; numerator and denominator equivalence alone are not
enough for empirical likelihood.  Unknown-order local power is only an event
transfer and requires order consistency under the local law \(P_{\beta_T}\).
\end{remark}

\section{Multiple-Break Extension Sketch}
\label{sec:multiple-break-sketch}

The main estimated-partition Wilks theorem is deliberately one-break.  For a
fixed number \(q_0>1\), the natural extension is a telescoping argument rather
than a new empirical-likelihood expansion.  For a local candidate
\(\Lambda=(k_1,\ldots,k_{q_0})\), define
\[
    \Lambda^{(0)}=\Lambda_0,
\]
and
\[
    \Lambda^{(j)}
    =(k_1,\ldots,k_j,k_{j+1,0},\ldots,k_{q_0,0}),
    \qquad j=1,\ldots,q_0.
\]
Then
\[
    \bm M_{K,\Lambda}-\bm M_{K,\Lambda_0}
    =
    \sum_{j=1}^{q_0}
    \{\bm M_{K,\Lambda^{(j)}}-
      \bm M_{K,\Lambda^{(j-1)}}\}.
\]
Each summand moves one boundary between two long adjacent segments.  If the
true breaks are separated and \(q_0\) is fixed, the one-break boundary,
coefficient-update, local-bias, and projection-noise arguments can be applied to
each summand, with only fixed multiplicative constants.  The local candidate set
has cardinality \(O(r_T^{q_0})\), so the uniform projection bounds replace
\(\log(1+r_T)\) by \(q_0\log(1+r_T)=O(\log r_T)\).

This sketch is not used in the main theorem.  A paper that claims general
fixed-\(q_0\) estimated-partition Wilks should promote this sketch to a full
lemma and verify the corresponding multi-boundary localization and local-law
order-selection statements.
\section{End-to-End Assumption Audit}
\label{sec:bp-end-to-end-audit}

This final audit records the rebuilt paper architecture.  The main theorem is
not a general martingale, general-multiple-break result.  It is a conditional-
on-design, sub-Gaussian, one-break estimated-partition Wilks theorem, with
unknown-order selection allowed over \(q=0,1,\ldots,\bar q\) when the true order
is \(q_0=1\).

\begin{center}
\small
\begin{tabular}{>{\raggedright\arraybackslash}p{0.25\textwidth}
                >{\raggedright\arraybackslash}p{0.42\textwidth}
                >{\raggedright\arraybackslash}p{0.23\textwidth}}
\hline
Proof block & Required inputs & What it delivers \\
\hline
Conditional-design environment & Conditional independence, mean zero, bounded variance, and sub-Gaussian tails given \(\D_T\) & Full-sample residualized scores and selected partitions are design-measurable, avoiding predictable-score assumptions \\
Known-partition Wilks & Variance convergence, nondegeneracy, score envelope, convex hull, and sieve feasibility & \(\ell_T(\beta_0,\Lambda_0)\Rightarrow\chi_1^2\) for fixed known \(q_0\) \\
One-break segmentation & Global Gram stability, deterministic profile identification, \(R_T=K+\log T+Ta_K^2\), and \(r_T=R_T/\Delta_T^2=o(\sqrt T)\) & \(|\widehat k-k_0|=O_p(r_T)\) under \(q_0=1\) \\
Estimated-break bridge & Boundary size, local bias stability, local full rank, search-uniform projection noise, and bridge rates & \(S_T(\widehat k)-S_T(k_0)=o_p(1)\) and \(V_T(\widehat k)-V_T(k_0)=o_p(1)\) \\
Unknown order & \(\operatorname{pen}_T(q,K)=q\kappa_TR_T\), \(\kappa_T\to\infty\), and \(\kappa_TR_T=o(T\Delta_T^2)\) & \(\widehat q(\beta_0)=1\) with probability tending to one \\
Local alternatives & \(\beta_T=\beta_0+h/\sqrt T\), \(\Gamma_T\to_p\Gamma\), \(\eta_T^2\to_p\eta^2\), and the local bridge & Noncentral limit \(\chi_1^2(h^2\Gamma^2/\eta^2)\) \\
Multiple breaks & Telescoping over boundaries plus uniform multi-boundary localization & Extension sketch only; not part of the main theorem \\
\hline
\end{tabular}
\end{center}

\paragraph{Audit conclusion.}
The rebuilt manuscript now has one probabilistic spine.  Conditional on
\(\D_T\), the EL score weights and the break search are fixed-design objects,
so the stochastic tools are conditional CLT, conditional LLN, and finite-search
sub-Gaussian concentration.  The price is explicit: the main estimated-break
Wilks and unknown-order Wilks results are stated for one nuisance break.  A
general fixed-\(q_0\) claim requires the telescoping extension in Section
\ref{sec:multiple-break-sketch} to be written as a full proof.

\section{Master Dependency Checklist}

The remaining obligations for a paper version are now concrete rather than
broad modular placeholders.

\begin{center}
\small
\begin{tabular}{>{\raggedright\arraybackslash}p{0.28\textwidth}
                >{\raggedright\arraybackslash}p{0.46\textwidth}
                >{\raggedright\arraybackslash}p{0.16\textwidth}}
\hline
Obligation & Concrete check & Status in this workbook \\
\hline
Design exogeneity & State and defend conditional independence/sub-Gaussianity given \(\D_T\) for the intended empirical or simulation design & Main assumption \\
Profile identification & Verify \(D_T(k)\geq c_\Delta |k-k_0|\Delta_T^2\) and the underfit gap from sieve jump energy and Gram stability & Isolated condition \\
Rate compatibility & Exhibit nonempty choices of \(K_T\), smoothness, and \(\Delta_T\) satisfying the bridge, projection, and penalty rates & Needs corollary \\
Convex hull & Give a primitive sign-mass condition or use an adjusted EL variant & Still separate EL regularity \\
Local drift & Estimate or assume \(\Gamma_T\to_p\Gamma\) and verify the quadratic/max drift bounds & Stated in local theorem \\
Multiple breaks & Promote the telescoping sketch to uniform multi-boundary lemmas & Future extension \\
\hline
\end{tabular}
\end{center}

\subsection*{Recommended remaining proof order}
\begin{enumerate}
    \item Add a rate-compatibility corollary, for example under
    \(\zeta_K^2\asymp K\), \(a_K\asymp K^{-s/d}\), and \(\Delta_T\geq c>0\).
    \item Add a primitive convex-hull lemma or switch to adjusted empirical
    likelihood.
    \item Verify profile identification for the chosen basis and design class.
    \item Decide whether the multiple-break telescoping sketch should remain an
    appendix note or become a theorem.
\end{enumerate}\end{document}
